%% 08-interconnection.tex -- generated by tools/build_from_docx.py; edit freely

\hypertarget{interconnection}{%
\section{Interconnection}\label{interconnection}}

The preceding four chapters described what each entity of the tuple
does. This chapter describes how they reach one another, and how an
agent reaches the platform that hosts it. Nothing in the framework works
if these paths are left to convention: an agent whose telemetry travels
by one route in the radio domain and another in transport is not
portable, and a control message whose delivery guarantee differs by
deployment is not a contract.

Two separations govern the chapter. The first is between control and
data: the service-based management bus carries control, and
observational data travels the Data Continuum. The second is between
what the platform provides and what the agent chooses. An agent declares
which flows it needs and with what properties; it does not select
transports, create topics or address peers. The platform binds declared
flows to concrete channels at activation, which is what allows the same
descriptor to be admitted in a domain whose substrate differs from the
one it was written against.

\hypertarget{the-control-plane}{%
\subsection{The Control Plane}\label{the-control-plane}}

Control among domains and the inter-domain management function travels
the cross-domain service-based management bus, realized in the
consortium's architecture over Apache Kafka. The bus carries
registration, lifecycle, instantiation, liveness and coordination
messages. It does not carry telemetry. This is a correction of the
earlier architecture description, in which observational data and
control shared the bus, and the reason for the correction is capacity
rather than tidiness: a bus sized for lifecycle messaging degrades under
measurement volume, and the degradation appears first as delayed control
at exactly the moment a service is in trouble.

\hypertarget{topic-hierarchy}{%
\subsubsection{Topic hierarchy}\label{topic-hierarchy}}

The bus organizes into four levels of topic, each with a distinct
lifetime.

A fixed pair of well-known registration topics, one carrying
registration requests and the other the responses, is the one entry
point by which a domain joins the system. A domain announces an
abstraction of its capabilities, its resources and its connections
toward neighboring transport domains; the request carries a temporary
identifier so that previously registered domains discard responses not
addressed to them, and the response carries the identifier the
inter-domain management function assigns together with the topics for
dedicated communication. Their lifetime is that of the system.

A pair of topics per domain carries bidirectional dedicated
communication: liveness heartbeats, resource updates and instantiation
messages, the last of which carry the details of the per-service or
per-agent topic. The pair is created by the inter-domain management
function when the domain registers. It is retained while the domain is
unreachable and is removed, along with its subscriptions, only when the
domain is deregistered, as Section 8.1.3 describes. Its lifetime is that
of the domain's membership.

A topic per service request is created by the inter-domain management
function before it dispatches the first instantiation command, and each
domain involved in realizing the request subscribes to it on receipt of
the command for its sub-service component. It carries coordination and
synchronization among the domains for the life of the service, including
termination, and the inter-domain management function deletes it, with
its subscriptions, once every participating domain has confirmed
deactivation. Its lifetime is that of the service instance.

A topic per agent instance is created by the orchestrator to which the
agent's descriptor is issued --- the inter-domain management function
for a global agent, the domain management function for a local one ---
before the first instantiation command, and each participating domain
subscribes to it on receipt of its portion of the service tree. It
carries the agent's instantiation status, binding status, activation
results, authority-state changes and assurance for the life of the
instance, and it disappears with its subscriptions when the instance is
terminated. It is distinct from the topic of any service, and an agent's
exchanges do not travel on a service topic.

The hierarchy matters to agents because it determines what an agent can
be addressed by. An agent is reachable within the domain that hosts it
and on the topic of its own instance; it is not addressable across the
system at large. A framework in which any agent could message any other
would have no containment story at the communication layer, whatever its
action-space declarations said.

\hypertarget{agents-do-not-manage-topics}{%
\subsubsection{Agents do not manage topics}\label{agents-do-not-manage-topics}}

Topic creation, subscription and teardown are performed by the
inter-domain management function and the domain management functions
through their service adapters. An agent declares the flows its
capabilities require, and the platform binds those flows to topics at
activation and releases them at deactivation. An agent that created its
own channel would outlive its own containment, because the channel would
survive a lifecycle action taken against the agent.

This is also what makes clean-up an obligation of the platform rather
than a matter of agent hygiene. Service and agent topics are
decommissioned when the lifecycle of the service or agent instance
expires, which prevents the resource leakage that accumulates when
teardown depends on well-behaved participants.

\hypertarget{unreachability-is-not-deregistration}{%
\subsubsection{Unreachability is not deregistration}\label{unreachability-is-not-deregistration}}

A domain whose heartbeat is lost beyond the configured tolerance is
unreachable. Unreachable is a condition of the domain and not a judgment
about the agents placed in it, and the distinction governs what the
platform is permitted to do next.

An unreachable domain's agents retain their registered authority. A
missed heartbeat is evidence about the path to the domain; it is
evidence about the agents only if one assumes the path and the agents
fail together, which is precisely the assumption that is false in the
common case. A platform that withdrew authority on unreachability would
strip healthy agents of it whenever the bus hiccuped, and would do so
most often in exactly the deployments --- testbeds, early integrations
--- where connectivity is least stable and where the resulting churn is
least affordable.

Retaining authority is safe because reachability already silences the
agent. An agent whose domain the platform cannot reach cannot reach the
conflict fabric either, and an intent that has not been mediated is
refused at the last mile. The agent is therefore unable to act for as
long as the condition persists, without any authority state being
changed to achieve it. The containment is a consequence of the
architecture rather than of a lifecycle operation, which is why no
lifecycle operation is warranted.

What the platform does do is report. The domain is reported as
unreachable rather than as silent, assurance records the condition
against every service with components in that domain, and objective
status reflects that those components are unverifiable for the duration.
A service whose assurance remained healthy while a participating domain
was unreachable would be asserting a service level nobody is currently
observing.

The agent resolves its own condition first. Each agent declares the
operational timeout within which it expects its own exchanges to
complete, set against the operations it performs, so that an agent whose
work legitimately takes long windows declares long ones. An agent that
passes its declared timeout without obtaining what it needs has
established that it is itself impaired, and reports itself unhealthy
rather than waiting to be diagnosed from outside. The platform does not
need external supervision to reach a conclusion the agent can reach
about itself, and an agent that can recover its own services does so
without any platform action.

The platform's recovery window is therefore the outer clock rather than
the only one, and is set wider than the agents' declared timeouts
because it already accounts for them. It governs the case the agent
cannot resolve: the agent's own resolution has failed, or the agent is
unreachable and cannot report anything at all. Uncertainty resolves in
every case and is never left standing. The window ends in one of two
outcomes and never in a continued unknown: the domain returns and its
agents are re-confirmed under Section 8.1.3.1, or the window expires and
deregistration proceeds.

The topic pair is retained for a recovery window that is longer than the
heartbeat tolerance. Deregistration, teardown of the pair and retirement
of the agents' generations occur together, and only when that window
expires. The two parameters are deployment values: the tolerance states
how long silence is treated as a transient, and the window states how
long the platform holds a place for a domain that may return.

\hypertarget{return}{%
\paragraph{Return}\label{return}}

On a heartbeat returning within the recovery window, the platform
re-confirms rather than re-gates. It checks the generation and
descriptor digest of each agent in the domain against the registry, and
where both are unchanged the agents resume under the authority they
already held.

A full activation gate run is required in three cases: the generation
has changed, the descriptor digest differs from the one registered, or
the agent missed evaluation cycles while unreachable such that its
conformance is no longer established. Each of these is a change in what
the agent is or in what is known about it, which is what the gate exists
to assess.

Re-gating in the absence of any of them costs the service an outage and
establishes nothing that was not already established. The purpose of the
gate is to prove that conditions hold before authority is granted; where
authority was never withdrawn and no condition changed, there is nothing
to re-prove.

\hypertarget{the-data-plane}{%
\subsection{The Data Plane}\label{the-data-plane}}

Observational data travels the Data Continuum rather than the management
bus. Measurements published by a Monitoring System capability, the
observations an Analytics Engine consumes, and the read-back evidence an
Actuator records are data products, and they are distributed by the
substrate the platform provides for that purpose rather than by the
control bus.

The framework deliberately does not name that substrate beyond its
properties. Work within the project on the unification of the monitoring
system is expected to define a broader capability for managing
monitoring across the Data Continuum and other sources, and a chapter
that hard-wired a product name would either constrain that work or be
overtaken by it. What the framework does fix are the properties a data
path must satisfy: measurements carry the observation envelope of
Chapter 4 unchanged, ordering is preserved within a metric and object,
and delivery latency is observable so that freshness can be evaluated
rather than assumed.

\hypertarget{flow-declarations}{%
\subsubsection{Flow declarations}\label{flow-declarations}}

Each flow is declared in the interfaces section of the descriptor with
the fields of Table 3-23: its kind, its producer and consumer, the
schema of its object, the binding classes permitted for it, its maximum
rate and its freshness, together with the delivery guarantee of Section
8.5 and, on a real-time loop, the latency bound it achieves. The
freshness, with the latency class the agent declares, states the timing
a channel has to meet. A flow does not name a transport, an address or a
topic. Channel selection is the responsibility of the domain's
orchestrator, the DMO: it reads the demands stated in the descriptor and
selects, from the binding classes the descriptor permits, a channel that
meets them. Where no available channel satisfies a declared flow, it
refuses activation rather than binding to a channel that satisfies the
declaration partially. A flow silently bound to a slower path than it
declared would produce an agent that appears conformant and misses its
latency class under load.

The binding classes are those of Section 3.5.10, and which may be
permitted follows from what the flow carries.
LOCAL\_\allowbreak{}RUNTIME is in-process, available where producer and
consumer are co-located. DATA\_\allowbreak{}CONTINUUM carries
observation streams and derived artifacts between instances. SMB carries
intents, outcomes and compact insight across a runtime boundary, as
control and coordination content only. SOUTHBOUND is the Actuator's
request-response interaction with its controller. An OBSERVATION flow is
never bound to the SMB, and an INTENT or OUTCOME flow is never bound to
the Data Continuum. Table 8-1 states the flow declaration fields and the
binding the DMO resolves.

\begin{wptable}{8-1}

\begin{longtable}[]{@{}
  >{\raggedright\arraybackslash}p{(\columnwidth - 6\tabcolsep) * \real{0.2625}}
  >{\raggedright\arraybackslash}p{(\columnwidth - 6\tabcolsep) * \real{0.1250}}
  >{\raggedright\arraybackslash}p{(\columnwidth - 6\tabcolsep) * \real{0.1000}}
  >{\raggedright\arraybackslash}p{(\columnwidth - 6\tabcolsep) * \real{0.5125}}@{}}
\caption{Flow declaration fields and the resolved
binding}\tabularnewline
\toprule\noalign{}
\begin{minipage}[b]{\linewidth}\raggedright
\textbf{Field in interfaces.\allowbreak{}flows{[}{]}}
\end{minipage} & \begin{minipage}[b]{\linewidth}\raggedright
\textbf{Type}
\end{minipage} & \begin{minipage}[b]{\linewidth}\raggedright
\textbf{Req.}
\end{minipage} & \begin{minipage}[b]{\linewidth}\raggedright
\textbf{Definition}
\end{minipage} \\
\midrule\noalign{}
\endfirsthead
\toprule\noalign{}
\begin{minipage}[b]{\linewidth}\raggedright
\textbf{Field in interfaces.\allowbreak{}flows{[}{]}}
\end{minipage} & \begin{minipage}[b]{\linewidth}\raggedright
\textbf{Type}
\end{minipage} & \begin{minipage}[b]{\linewidth}\raggedright
\textbf{Req.}
\end{minipage} & \begin{minipage}[b]{\linewidth}\raggedright
\textbf{Definition}
\end{minipage} \\
\midrule\noalign{}
\endhead
\bottomrule\noalign{}
\endlastfoot
\textbf{flow} & enum & M & OBSERVATION, INSIGHT, INTENT or OUTCOME.
Determines the binding classes that may be permitted. \\
\textbf{from, to} & enum & M & Producing and consuming capability, or a
provider or consumer outside the agent. \\
\textbf{schema} & schema-ref & M & Schema of the flow's object and its
version; compatibility is checked at activation (Section 8.6). \\
\textbf{binding\_\allowbreak{}classes} & list of enum & M & Binding
classes permitted for the flow, within those allowed for its kind. The
descriptor permits; it never selects. \\
\textbf{max\_\allowbreak{}rate} & rate & M & Maximum rate at which the
producer emits. \\
\textbf{freshness} & duration & M & Maximum age at delivery; with the
agent's latency class, the timing a channel has to meet. \\
\textbf{delivery\_\allowbreak{}guarantee} & enum & M &
at\_\allowbreak{}most\_\allowbreak{}once,
at\_\allowbreak{}least\_\allowbreak{}once or exactly\_\allowbreak{}once
(Section 8.5). \\
\textbf{latency\_\allowbreak{}bound} & object & C & On a real-time loop,
the latency the flow achieves and its evidence (Section 3.8.3). \\
\textbf{Resolved binding} & enum & --- & LOCAL\_\allowbreak{}RUNTIME,
DATA\_\allowbreak{}CONTINUUM, SMB or SOUTHBOUND, selected by the DMO
from the permitted classes and recorded in the resolved configuration;
never selected by the agent. \\
\end{longtable}

\end{wptable}

\hypertarget{paths-between-the-roles}{%
\subsection{Paths Between the Roles}\label{paths-between-the-roles}}

The four roles of the tuple exchange along fixed paths, and the fixity
is what allows an agent's internal composition to vary without its
interfaces varying.

A Monitoring System capability publishes observation envelopes to the
Analytics Engine capabilities within its agent and, where its descriptor
declares the flow, to the domain management function. An Analytics
Engine capability publishes insight envelopes to the Decision Engine
capabilities of its agent, and knowledge crosses domain boundaries
freely as described in Chapter 5. A Decision Engine capability emits
intent envelopes through the conflict fabric binding of Chapter 6, and
never directly to an Actuator. An Actuator publishes execution envelopes
to the originating Decision Engine and to the lifecycle and evaluation
path of the domain management function, dual-routed by the bus rather
than addressed by the agent.

Two of these paths are mediated rather than direct, and the mediation is
the substance of the containment described in the preceding chapters.
Where a path appears direct in a diagram and mediated in the text, the
text governs.

\hypertarget{cross-agent-exchange}{%
\subsubsection{Cross-agent exchange}\label{cross-agent-exchange}}

Agents within a domain exchange insights freely. Decisions may be
exchanged laterally as proposals, but a lateral proposal binds no one
and confers no authority; the final decision rests with the overlooking
agent, which is the inter-domain management function. Authority descends
and never passes sideways. Execution feedback is not exchanged between
agents at all, because an actuation outcome is evidence about the
network held by the platform rather than a message between participants.

\hypertarget{external-interfaces}{%
\subsection{External Interfaces}\label{external-interfaces}}

An agent reaches the network only through the controller bindings of
Chapter 7 and reaches the management system only through its declared
flows. It holds no other external interface. A capability that opened a
direct connection to a device, a database or an external service would
be outside every enforcement point the framework defines, and its
declaration would describe an agent that is not the one running.

Where a capability genuinely requires an external source, that source is
declared in the descriptor as a source of the appropriate kind and is
reached through the platform, which applies the trust-boundary handling
of Chapter 4 and records the dependency for evaluation.

\hypertarget{delivery-and-failure-semantics}{%
\subsection{Delivery and Failure Semantics}\label{delivery-and-failure-semantics}}

A declared delivery guarantee is a property the platform enforces and
not an aspiration the agent holds. Table 8-2 states the three available
guarantees and their permitted uses. At most once suits high-rate
measurement where a lost sample is preferable to a delayed one. At least
once suits control messaging whose handler is idempotent. Exactly once
is available on the management bus for lifecycle messaging where
duplicate application would be damaging, and it is not permitted for
OBSERVATION flows, at whose rates its cost is not repaid.

The interaction between delivery semantics and the actuation outcomes of
Chapter 7 is worth stating plainly, because it is the place where a
messaging property becomes a safety property. An at-least-once control
flow may deliver an instantiation or intent message twice; the
idempotency declaration of the receiving capability is what makes that
safe, and a non-idempotent action type reached by an at-least-once flow
is a declaration error that admission rejects.

Where a flow's channel fails, the capability's declared failure behavior
applies as described in the chapter for its role, and the platform
raises a flow fault against the binding rather than against the agent's
logic. This distinction matters for evaluation: an agent starved of
input by a broken channel is not an agent behaving badly, and conflating
the two would penalize participants for infrastructure they do not
control.

\begin{wptable}{8-2}

\begin{longtable}[]{@{}
  >{\raggedright\arraybackslash}p{(\columnwidth - 4\tabcolsep) * \real{0.2000}}
  >{\raggedright\arraybackslash}p{(\columnwidth - 4\tabcolsep) * \real{0.4000}}
  >{\raggedright\arraybackslash}p{(\columnwidth - 4\tabcolsep) * \real{0.4000}}@{}}
\caption{Delivery guarantees and their permitted uses}\tabularnewline
\toprule\noalign{}
\begin{minipage}[b]{\linewidth}\raggedright
\textbf{Guarantee}
\end{minipage} & \begin{minipage}[b]{\linewidth}\raggedright
\textbf{Intended use}
\end{minipage} & \begin{minipage}[b]{\linewidth}\raggedright
\textbf{Constraint}
\end{minipage} \\
\midrule\noalign{}
\endfirsthead
\toprule\noalign{}
\begin{minipage}[b]{\linewidth}\raggedright
\textbf{Guarantee}
\end{minipage} & \begin{minipage}[b]{\linewidth}\raggedright
\textbf{Intended use}
\end{minipage} & \begin{minipage}[b]{\linewidth}\raggedright
\textbf{Constraint}
\end{minipage} \\
\midrule\noalign{}
\endhead
\bottomrule\noalign{}
\endlastfoot
\textbf{at\_\allowbreak{}most\_\allowbreak{}once} & High-rate
observational flows where a lost sample is preferable to a delayed one.
& Loss is reported through the coverage and loss signals of Chapter 4,
never as zero. \\
\textbf{at\_\allowbreak{}least\_\allowbreak{}once} & Control and
lifecycle messaging whose handler is idempotent. & A non-idempotent
action type reached by an at-least-once flow is rejected at
admission. \\
\textbf{exactly\_\allowbreak{}once} & Lifecycle messaging on the
management bus where duplicate application would be damaging. & Not
available for high-rate observational flows. \\
\end{longtable}

\end{wptable}

\hypertarget{schema-versioning}{%
\subsection{Schema Versioning}\label{schema-versioning}}

Every flow declares a payload schema and its version, and the platform
checks compatibility at activation rather than discovering
incompatibility in traffic. A change that adds an optional field, widens
a range or adds an enumeration value a consumer may ignore is
compatible, and existing bindings survive it. A change that removes a
field, narrows a range, changes a unit or changes the meaning of an
existing value is incompatible, and it requires a new schema version and
a rebinding of every affected flow.

A consumer bound to an incompatible producer version is not activated,
and the refusal is reported to the requester of the instantiation as a
diagnostic naming the flow, the two schema versions and the specific
incompatibility between them. A refusal that does not say why is of no
use to the partner holding the descriptor, who is the only party able to
correct it. The alternative, in which the platform binds and the
mismatch surfaces as malformed payloads at runtime, produces failures
attributed to the wrong component, which is how a versioning problem
becomes a week of misdirected debugging.

\hypertarget{requirements-5}{%
\subsection{Requirements}\label{requirements-5}}

\begin{requirement}{Req-T23-IC-01}{MUST}{Control and data are separated}
\reqfield{Description / rationale}{%
The service-based management bus shall carry control, lifecycle,
registration, liveness and coordination messaging only. Observational
data shall travel the Data Continuum. An OBSERVATION flow shall never be
bound to the SMB, and an INTENT or OUTCOME flow shall never be bound to
the Data Continuum.}
\reqfield{Dependency}{%
Req-T23-FW-02}
\reqfield{Traceability (Forward)}{%
\begin{itemize}
\tightlist
\item
  Binding kind checked by the platform at activation.
\end{itemize}}
\reqfield{UNITY-6G components}{%
\begin{itemize}
\tightlist
\item
  SMB
\item
  Data Continuum
\item
  Platform runtime
\end{itemize}}
\reqfield{Linked KPIs}{%
\begin{itemize}
\tightlist
\item
  KPI\#24: Decrease end-to-end message delivery latency and increase
  maximum message publish throughput in the SBA service bus in
  integrated network scenarios.
\item
  KPI\#4: Real-time reconfiguration of IEEE 802.11 network domain in
  range smaller than 10 ms from taking the decision.
\end{itemize}}
\end{requirement}

\begin{requirement}{Req-T23-IC-02}{MUST}{Observation envelopes travel unchanged}
\reqfield{Description / rationale}{%
A measurement carried on a data path shall carry the observation
envelope of Chapter 4 unchanged, ordering shall be preserved within a
metric and object, and delivery latency shall be observable so that
freshness is evaluated rather than assumed.}
\reqfield{Dependency}{%
Req-T23-MS-01}
\reqfield{Traceability (Forward)}{%
\begin{itemize}
\tightlist
\item
  Envelope integrity checked at consumption.
\end{itemize}}
\reqfield{UNITY-6G components}{%
\begin{itemize}
\tightlist
\item
  Data Continuum
\item
  MS runtime
\item
  AE runtime
\end{itemize}}
\reqfield{Linked KPIs}{%
\begin{itemize}
\tightlist
\item
  KPI\#36: Guarantee 99.99999\% service availability and reliability
  based on effective monitoring and LCM.
\item
  KPI\#1: Deliver four distributed management entities (MS, AE, DE and
  ACT) deployable in each heterogeneous domain.
\end{itemize}}
\end{requirement}

\begin{requirement}{Req-T23-IC-03}{MUST}{Flows are declared, transports are not}
\reqfield{Description / rationale}{%
Each flow shall be declared with the fields of Table 3-23, including its
delivery guarantee. A flow shall not name a transport, an address or a
topic. The DMO shall select the binding of each flow from the binding
classes the descriptor permits and bind it at activation.}
\reqfield{Dependency}{%
Req-T23-FW-04}
\reqfield{Traceability (Forward)}{%
\begin{itemize}
\tightlist
\item
  Declaration validated at admission.
\item
  Binding performed by the platform at activation.
\end{itemize}}
\reqfield{UNITY-6G components}{%
\begin{itemize}
\tightlist
\item
  Tuple descriptor
\item
  Platform runtime
\end{itemize}}
\reqfield{Linked KPIs}{%
\begin{itemize}
\tightlist
\item
  KPI\#41: Reduction of 70\% of deployment time of AI-based
  resource/energy management algorithms.
\item
  KPI\#1: Deliver four distributed management entities (MS, AE, DE and
  ACT) deployable in each heterogeneous domain.
\item
  KPI\#4: Real-time reconfiguration of IEEE 802.11 network domain in
  range smaller than 10 ms from taking the decision.
\end{itemize}}
\end{requirement}

\begin{requirement}{Req-T23-IC-04}{MUST}{Partial satisfaction is not a binding}
\reqfield{Description / rationale}{%
Where no available channel satisfies a declared flow, the platform shall
refuse activation rather than bind to a channel that satisfies the
declaration partially.}
\reqfield{Dependency}{%
Req-T23-IC-03}
\reqfield{Traceability (Forward)}{%
\begin{itemize}
\tightlist
\item
  Enforced at activation by the platform.
\end{itemize}}
\reqfield{UNITY-6G components}{%
\begin{itemize}
\tightlist
\item
  Platform runtime
\item
  DMO
\end{itemize}}
\reqfield{Linked KPIs}{%
\begin{itemize}
\tightlist
\item
  KPI\#4: Real-time reconfiguration of IEEE 802.11 network domain in
  range smaller than 10 ms from taking the decision.
\item
  KPI\#36: Guarantee 99.99999\% service availability and reliability
  based on effective monitoring and LCM.
\end{itemize}}
\end{requirement}

\begin{requirement}{Req-T23-IC-05}{MUST}{Fixed paths between the roles}
\reqfield{Description / rationale}{%
MS capabilities shall publish observation envelopes to the AE
capabilities of their agent and, where declared, to the DMO. AE
capabilities shall publish insight envelopes to the DE capabilities of
their agent. DE capabilities shall emit intent through the conflict
fabric binding and never directly to an ACT. ACT capabilities shall
publish execution envelopes to the originating DE and to the DMO
lifecycle and evaluation path.}
\reqfield{Dependency}{%
Req-T23-FW-02}
\reqfield{Traceability (Forward)}{%
\begin{itemize}
\tightlist
\item
  Paths established by the platform at activation.
\end{itemize}}
\reqfield{UNITY-6G components}{%
\begin{itemize}
\tightlist
\item
  Platform runtime
\end{itemize}}
\reqfield{Linked KPIs}{%
\begin{itemize}
\tightlist
\item
  KPI\#1: Deliver four distributed management entities (MS, AE, DE and
  ACT) deployable in each heterogeneous domain.
\item
  KPI\#21: Reduce the performance penalty due to conflicts by 20\% by
  means of DLT, hierarchical algorithms and policy-based conflict
  resolution schemes.
\end{itemize}}
\end{requirement}

\begin{requirement}{Req-T23-IC-06}{MUST}{Lateral exchange confers no authority}
\reqfield{Description / rationale}{%
Agents may exchange insights freely and may exchange decisions laterally
as proposals. A lateral proposal shall bind no recipient and shall
confer no authority; the final decision shall rest with the overlooking
agent. Execution feedback shall not be exchanged between agents.}
\reqfield{Dependency}{%
Req-T23-DE-11}
\reqfield{Traceability (Forward)}{%
\begin{itemize}
\tightlist
\item
  Authority checked at the conflict fabric and the IDMO.
\end{itemize}}
\reqfield{UNITY-6G components}{%
\begin{itemize}
\tightlist
\item
  IDMO
\item
  Conflict fabric
\item
  SMB
\end{itemize}}
\reqfield{Linked KPIs}{%
\begin{itemize}
\tightlist
\item
  KPI\#21: Reduce the performance penalty due to conflicts by 20\% by
  means of DLT, hierarchical algorithms and policy-based conflict
  resolution schemes.
\item
  KPI\#26: Increase the ratio of service management and LCM tasks
  resolved by local AE/DE components.
\item
  KPI\#1: Deliver four distributed management entities (MS, AE, DE and
  ACT) deployable in each heterogeneous domain.
\end{itemize}}
\end{requirement}

\begin{requirement}{Req-T23-IC-07}{MUST}{Topics are managed by the platform}
\reqfield{Description / rationale}{%
Topic creation, subscription and teardown shall be performed by the IDMO
and DMO through their service adapters. An agent shall not create,
subscribe to or tear down a topic. Service topics shall be
decommissioned when the service lifecycle expires.}
\reqfield{Dependency}{%
Req-T23-IC-03}
\reqfield{Traceability (Forward)}{%
\begin{itemize}
\tightlist
\item
  Service adapter is the only creation path.
\item
  Teardown driven by the service lifecycle.
\end{itemize}}
\reqfield{UNITY-6G components}{%
\begin{itemize}
\tightlist
\item
  IDMO
\item
  DMO
\item
  SMB
\end{itemize}}
\reqfield{Linked KPIs}{%
\begin{itemize}
\tightlist
\item
  KPI\#24: Decrease end-to-end message delivery latency and increase
  maximum message publish throughput in the SBA service bus in
  integrated network scenarios.
\item
  KPI\#31: Reduce service recovery time below 180 s.
\item
  KPI\#26: Increase the ratio of service management and LCM tasks
  resolved by local AE/DE components.
\end{itemize}}
\end{requirement}

\begin{requirement}{Req-T23-IC-08}{MUST}{No external interface outside declared flows}
\reqfield{Description / rationale}{%
A capability shall reach the network only through the controller
bindings of Chapter 7 and the management system only through its
declared flows. An external source required by a capability shall be
declared in the descriptor and reached through the platform, which shall
apply trust-boundary handling and record the dependency for evaluation.}
\reqfield{Dependency}{%
Req-T23-IC-03}
\reqfield{Traceability (Forward)}{%
\begin{itemize}
\tightlist
\item
  Checked at admission.
\item
  Undeclared egress detected by the platform runtime.
\end{itemize}}
\reqfield{UNITY-6G components}{%
\begin{itemize}
\tightlist
\item
  Platform runtime
\item
  Admission (IDMO/DMO)
\end{itemize}}
\reqfield{Linked KPIs}{%
\begin{itemize}
\tightlist
\item
  KPI\#41: Reduction of 70\% of deployment time of AI-based
  resource/energy management algorithms.
\item
  KPI\#36: Guarantee 99.99999\% service availability and reliability
  based on effective monitoring and LCM.
\item
  KPI\#31: Reduce service recovery time below 180 s.
\end{itemize}}
\end{requirement}

\begin{requirement}{Req-T23-IC-09}{MUST}{Delivery guarantees are enforced and constrained}
\reqfield{Description / rationale}{%
The platform shall enforce the declared delivery guarantee.
exactly\_\allowbreak{}once shall not be permitted for OBSERVATION flows,
and an action type declared non-idempotent shall not be reached by an
at\_\allowbreak{}least\_\allowbreak{}once flow; such a declaration shall
be rejected at admission.}
\reqfield{Dependency}{%
Req-T23-ACT-05}
\reqfield{Traceability (Forward)}{%
\begin{itemize}
\tightlist
\item
  Combination checked at admission against the action contract.
\end{itemize}}
\reqfield{UNITY-6G components}{%
\begin{itemize}
\tightlist
\item
  Platform runtime
\item
  Admission (IDMO/DMO)
\end{itemize}}
\reqfield{Linked KPIs}{%
\begin{itemize}
\tightlist
\item
  KPI\#36: Guarantee 99.99999\% service availability and reliability
  based on effective monitoring and LCM.
\item
  KPI\#31: Reduce service recovery time below 180 s.
\end{itemize}}
\end{requirement}

\begin{requirement}{Req-T23-IC-10}{MUST}{Schema compatibility is checked at activation}
\reqfield{Description / rationale}{%
The platform shall check payload schema compatibility at activation.
Adding an optional field, widening a range or adding an ignorable
enumeration value shall be compatible; removing a field, narrowing a
range, changing a unit or changing the meaning of a value shall be
incompatible and shall require a new version and rebinding. A consumer
bound to an incompatible producer version shall not be activated, and
the platform shall report the refusal to the requester of the
instantiation, naming the flow, both schema versions and the
incompatibility.}
\reqfield{Dependency}{%
Req-T23-IC-03}
\reqfield{Traceability (Forward)}{%
\begin{itemize}
\tightlist
\item
  Compatibility evaluated at activation, not in traffic.
\item
  Refusal carries a diagnostic to the requester.
\end{itemize}}
\reqfield{UNITY-6G components}{%
\begin{itemize}
\tightlist
\item
  Platform runtime
\end{itemize}}
\reqfield{Linked KPIs}{%
\begin{itemize}
\tightlist
\item
  KPI\#36: Guarantee 99.99999\% service availability and reliability
  based on effective monitoring and LCM.
\item
  KPI\#41: Reduction of 70\% of deployment time of AI-based
  resource/energy management algorithms.
\end{itemize}}
\end{requirement}

\hypertarget{contribution-to-the-project-kpis-4}{%
\subsection{CONTRIBUTION TO THE PROJECT KPIs}\label{contribution-to-the-project-kpis-4}}

Table 8-3 states what interconnection contributes to each KPI it bears
on, and Table 8-4 maps the requirements of this chapter onto those KPIs.

\begin{wptable}{8-3}

\begin{longtable}[]{@{}
  >{\raggedright\arraybackslash}p{(\columnwidth - 4\tabcolsep) * \real{0.1499}}
  >{\raggedright\arraybackslash}p{(\columnwidth - 4\tabcolsep) * \real{0.6375}}
  >{\raggedright\arraybackslash}p{(\columnwidth - 4\tabcolsep) * \real{0.2126}}@{}}
\caption{Contribution of interconnection to the project
KPIs}\tabularnewline
\toprule\noalign{}
\begin{minipage}[b]{\linewidth}\raggedright
\textbf{KPI}
\end{minipage} & \begin{minipage}[b]{\linewidth}\raggedright
\textbf{Contribution of interconnection}
\end{minipage} & \begin{minipage}[b]{\linewidth}\raggedright
\textbf{Requirements}
\end{minipage} \\
\midrule\noalign{}
\endfirsthead
\toprule\noalign{}
\begin{minipage}[b]{\linewidth}\raggedright
\textbf{KPI}
\end{minipage} & \begin{minipage}[b]{\linewidth}\raggedright
\textbf{Contribution of interconnection}
\end{minipage} & \begin{minipage}[b]{\linewidth}\raggedright
\textbf{Requirements}
\end{minipage} \\
\midrule\noalign{}
\endhead
\bottomrule\noalign{}
\endlastfoot
\textbf{KPI\#1} & Fixed paths between the four roles are what let an
agent\textquotesingle s internal composition vary while its interfaces
do not. & Req-T23-IC-05, Req-T23-IC-06 \\
\textbf{KPI\#4} & Latency classes are declared per flow and bound at
activation, so a sub-10 ms path is refused rather than silently bound to
a slower channel. & Req-T23-IC-03, Req-T23-IC-04 \\
\textbf{KPI\#24} & Keeping observational data off the management bus
leaves the bus to control messaging, and platform-managed topics are
torn down with the services they serve, which keeps the bus's delivery
latency and publish throughput within bounds as services multiply. &
Req-T23-IC-01, Req-T23-IC-07 \\
\textbf{KPI\#26} & Per-domain, per-service and per-agent topic scoping
is what allows a task to be resolved within the domain that holds it. &
Req-T23-IC-07 \\
\textbf{KPI\#31} & Platform-managed topic teardown at service
decommissioning prevents the resource leakage that slows recovery. &
Req-T23-IC-07, Req-T23-IC-08 \\
\textbf{KPI\#36} & Measurements travel with their observation envelope
intact and their delivery latency observable, flow faults are attributed
to the binding rather than the agent, and schema incompatibility blocks
activation rather than surfacing as runtime corruption. & Req-T23-IC-02,
Req-T23-IC-09, Req-T23-IC-10 \\
\textbf{KPI\#41} & An agent declaring flows rather than transports is
admissible in a domain whose substrate differs from the one it was
written against. & Req-T23-IC-03, Req-T23-IC-08 \\
\end{longtable}

\end{wptable}

\begin{wptable}{8-4}\footnotesize

\begin{longtable}[]{@{}
  >{\raggedright\arraybackslash}p{(\columnwidth - 16\tabcolsep) * \real{0.1813}}
  >{\raggedright\arraybackslash}p{(\columnwidth - 16\tabcolsep) * \real{0.1023}}
  >{\raggedright\arraybackslash}p{(\columnwidth - 16\tabcolsep) * \real{0.1023}}
  >{\raggedright\arraybackslash}p{(\columnwidth - 16\tabcolsep) * \real{0.1023}}
  >{\raggedright\arraybackslash}p{(\columnwidth - 16\tabcolsep) * \real{0.1023}}
  >{\raggedright\arraybackslash}p{(\columnwidth - 16\tabcolsep) * \real{0.1023}}
  >{\raggedright\arraybackslash}p{(\columnwidth - 16\tabcolsep) * \real{0.1023}}
  >{\raggedright\arraybackslash}p{(\columnwidth - 16\tabcolsep) * \real{0.1023}}
  >{\raggedright\arraybackslash}p{(\columnwidth - 16\tabcolsep) * \real{0.1025}}@{}}
\caption{Mapping of interconnection requirements onto the project
KPIs}\tabularnewline
\toprule\noalign{}
\begin{minipage}[b]{\linewidth}\raggedright
\textbf{KPI\textbackslash{}}

\textbf{Requirement}
\end{minipage} & \begin{minipage}[b]{\linewidth}\raggedright
\textbf{\#1}
\end{minipage} & \begin{minipage}[b]{\linewidth}\raggedright
\textbf{\#4}
\end{minipage} & \begin{minipage}[b]{\linewidth}\raggedright
\textbf{\#21}
\end{minipage} & \begin{minipage}[b]{\linewidth}\raggedright
\textbf{\#24}
\end{minipage} & \begin{minipage}[b]{\linewidth}\raggedright
\textbf{\#26}
\end{minipage} & \begin{minipage}[b]{\linewidth}\raggedright
\textbf{\#31}
\end{minipage} & \begin{minipage}[b]{\linewidth}\raggedright
\textbf{\#36}
\end{minipage} & \begin{minipage}[b]{\linewidth}\raggedright
\textbf{\#41}
\end{minipage} \\
\midrule\noalign{}
\endfirsthead
\toprule\noalign{}
\begin{minipage}[b]{\linewidth}\raggedright
\textbf{KPI\textbackslash{}}

\textbf{Requirement}
\end{minipage} & \begin{minipage}[b]{\linewidth}\raggedright
\textbf{\#1}
\end{minipage} & \begin{minipage}[b]{\linewidth}\raggedright
\textbf{\#4}
\end{minipage} & \begin{minipage}[b]{\linewidth}\raggedright
\textbf{\#21}
\end{minipage} & \begin{minipage}[b]{\linewidth}\raggedright
\textbf{\#24}
\end{minipage} & \begin{minipage}[b]{\linewidth}\raggedright
\textbf{\#26}
\end{minipage} & \begin{minipage}[b]{\linewidth}\raggedright
\textbf{\#31}
\end{minipage} & \begin{minipage}[b]{\linewidth}\raggedright
\textbf{\#36}
\end{minipage} & \begin{minipage}[b]{\linewidth}\raggedright
\textbf{\#41}
\end{minipage} \\
\midrule\noalign{}
\endhead
\bottomrule\noalign{}
\endlastfoot
\textbf{Req-T23-IC-01} & & \kpimark{} & & \kpimark{} & & & & \\
\textbf{Req-T23-IC-02} & \kpimark{} & & & & & & \kpimark{} & \\
\textbf{Req-T23-IC-03} & \kpimark{} & \kpimark{} & & & & & &
\kpimark{} \\
\textbf{Req-T23-IC-04} & & \kpimark{} & & & & & \kpimark{} & \\
\textbf{Req-T23-IC-05} & \kpimark{} & & \kpimark{} & & & & & \\
\textbf{Req-T23-IC-06} & \kpimark{} & & \kpimark{} & & \kpimark{} & &
& \\
\textbf{Req-T23-IC-07} & & & & \kpimark{} & \kpimark{} & \kpimark{} &
& \\
\textbf{Req-T23-IC-08} & & & & & & \kpimark{} & \kpimark{} &
\kpimark{} \\
\textbf{Req-T23-IC-09} & & & & & & \kpimark{} & \kpimark{} & \\
\textbf{Req-T23-IC-10} & & & & & & & \kpimark{} & \kpimark{} \\
\end{longtable}

\end{wptable}
